\begin{table}[ht]
\centering
\small
\begin{tabular}{lcccccc}
\toprule
\textbf{Retrieval System} & \textbf{Decisions} & \multicolumn{2}{c}{\textbf{Top-3 Occlusion (ROAR)}} & \multicolumn{2}{c}{\textbf{Outcome Shuffling}} & \textbf{Random Replacement} \\
\cmidrule(lr){3-4} \cmidrule(lr){5-6} \cmidrule(lr){7-7}
& & \textbf{Mean Shift} & \textbf{Rank $\rho$} & \textbf{Mean Shift} & \textbf{Rank $\rho$} & \textbf{Rank $\rho$ (Noise)} \\
\midrule
P0 (Global Learned Memory) & 335 & 211 bps & 0.801 & 251 bps & 0.693 & \textbf{0.004} \\
P0* (Domestic Guardrail) & 330 & 177 bps & 0.798 & 197 bps & 0.736 & \textbf{0.058} \\
P2 (Mean-Only Retrieval) & 335 & 199 bps & 0.841 & 233 bps & 0.750 & \textbf{-0.085} \\
P3 (Raw $k$NN Memory) & 342 & 174 bps & 0.840 & 214 bps & 0.764 & \textbf{-0.003} \\
\bottomrule
\end{tabular}
\caption{Interpretability Faithfulness Evaluation across 1,342 decision instances. Testing whether retrieved historical analogues causally govern portfolio selection rather than serving as decorative post-hoc explanations. Injecting uninformative historical noise (Random Replacement) collapses decision rank correlation to $\rho \approx 0.00$. Removing the top-3 highest-weighted precedents (Top-3 Occlusion) shifts expected return by 174--211 bps, demonstrating that the decision engine is causally anchored in the retrieved historical evidence.}
\label{tab:interpretability_faithfulness}
\end{table}
